\documentclass[10pt]{article}

% =========================
% Fonts & Encoding
% =========================
\usepackage[utf8]{inputenc}
\usepackage[T1]{fontenc}
\usepackage[sc]{mathpazo}      % Palatino w/ small caps

% =========================
% Layout
% =========================
\usepackage[a4paper, margin=1in]{geometry}
\usepackage{setspace}
\onehalfspacing
\setlength{\parindent}{0pt}
\setlength{\parskip}{1em}

% =========================
% Floats / placement
% =========================
\usepackage{float}
\usepackage{placeins}
\usepackage{subcaption}        % subfigure grids

% =========================
% Math & Tables
% =========================
\usepackage{amsmath,amsfonts}
\usepackage{booktabs}
\usepackage{array}
\usepackage{multirow}
\usepackage[table]{xcolor}

% =========================
% Figures & formatting
% =========================
\usepackage{graphicx}
\usepackage{microtype}
\usepackage{caption}
\usepackage{titlesec}
\usepackage{enumitem}
\usepackage{fancyhdr}

% =========================
% Unicode helpers
% =========================
\usepackage{newunicodechar}
\newunicodechar{Δ}{\ensuremath{\Delta}}
\newunicodechar{−}{\ensuremath{-}}
\newunicodechar{–}{--}
\newunicodechar{—}{---}
\newunicodechar{≤}{\ensuremath{\le}}
\newunicodechar{≥}{\ensuremath{\ge}}
\newunicodechar{≈}{\ensuremath{\approx}}
\newunicodechar{ρ}{\ensuremath{\rho}}

\errorcontextlines=20

% =========================
% Citations
% =========================
\usepackage[authoryear,round]{natbib}
\setcitestyle{aysep={,},citesep={;},open={(},close={)}}
\setcitestyle{maxcitenames=2,maxbibnames=99}

% =========================
% Hyperlinks (load late)
% =========================
\usepackage[colorlinks=true,
            linkcolor=blue,
            citecolor=blue,
            urlcolor=blue]{hyperref}

% =========================
% Section formatting
% =========================
\titleformat{\section}{\large\bfseries}{\thesection}{1em}{}
\titleformat{\subsection}{\normalsize\bfseries}{\thesubsection}{1em}{}

% =========================
% Header & footer
% =========================
\pagestyle{fancy}
\fancyhf{}
\rhead{Daniel Tesfai Kebede}
\lhead{Layer-Wise Gender Bias in Multilingual Transformers}
\cfoot{\thepage}

% =========================
% Metadata
% =========================
\newcommand{\faculty}{Faculty of Computational Linguistics and\\Digital Humanities}
\newcommand{\papertitle}{Layer-Wise Gender Bias in Multilingual Foundation Models:\\A Five-Architecture Cross-Lingual Analysis}
\newcommand{\authorname}{Daniel Tesfai Kebede}
\newcommand{\matrnr}{1716694}
\newcommand{\course}{Trends in Machine Learning}
\newcommand{\instructor}{[Instructor Name]}
\newcommand{\submissiondate}{[Submission Date]}

\begin{document}

% -------------------------------------------------------
%  Page 1 — Title Page
% -------------------------------------------------------
\begin{titlepage}
  \thispagestyle{empty}

  \noindent
  \begin{minipage}[t]{0.5\textwidth}
    \vspace*{0pt}%
    % \includegraphics[height=1.9cm]{media/Logo-Uni-Trier.jpg}
    {\small University of Trier}
  \end{minipage}%
  \begin{minipage}[t]{0.5\textwidth}
    \vspace*{0pt}\raggedleft
    {\small \faculty}
  \end{minipage}

  \vfill
  \begin{center}
    {\Large\bfseries \papertitle\par}
    \vspace{1em}
    {\normalsize Seminar Paper\par}
    \vspace{1.5em}
    {\normalsize submitted by\par}
    \vspace{0.6em}
    {\normalsize \textbf{\authorname} (Matr.\ \matrnr)\par}
  \end{center}
  \vfill

  {\normalsize
   \textbf{Course:}\quad \course\\[0.3em]
   \textbf{Submission:}\quad \submissiondate\\[0.3em]
   \textbf{Instructor:}\quad \instructor
  }
\end{titlepage}

% -------------------------------------------------------
%  Page 2 — Table of Contents
% -------------------------------------------------------
\newpage
\begingroup
  \thispagestyle{empty}
  \singlespacing
  \setlength{\parskip}{0pt}
  \setcounter{tocdepth}{2}
  \tableofcontents
\endgroup

% -------------------------------------------------------
%  Main Content
% -------------------------------------------------------
\newpage
\pagenumbering{arabic}
\setcounter{page}{1}

% =======================================================
\begin{abstract}
% =======================================================

Multilingual transformer models embed dozens of languages within shared
representational spaces, enabling cross-lingual transfer while raising
questions about how social associations are distributed across languages
with unequal training data. Although gender bias in English embeddings has
been extensively studied, comparatively little work examines how multilingual
contextual models encode such associations across both high-resource and
genuinely low-resource languages, or how those associations evolve across
model depth.

This study investigates gender bias in Tigrigna—a low-resource Ge'ez-script
language—as a primary case study, with comparative analysis across Spanish
and Arabic, using five multilingual transformer architectures:
\texttt{bert-base-multilingual-cased} (mBERT),
\texttt{xlm-roberta-base}, \texttt{xlm-roberta-large},
\texttt{facebook/xlm-v-base}, and \texttt{microsoft/mdeberta-v3-base}.
Contextual token representations are reduced to stable word-level embeddings
via contextual-to-static aggregation \citep{bommasani2020reductions}, enabling
layer-wise projection of profession terms onto a geometric gender direction
\citep{bolukbasi2016man}.

Our results reveal four key findings. First, four of five models encode
measurable male-leaning bias for Tigrigna despite the language's absence
from dedicated pretraining data. Second, bias accumulates progressively
across transformer depth, peaking at layer~9 (base models) or layer~23
(XLM-R-large). Third, two independent geometric estimators agree at
$\rho \geq 0.970$ across all valid models, confirming methodological
robustness. Fourth, and most strikingly, mDeBERTa-v3-base—trained on
comparable multilingual data—produces a complete directional reversal:
28 of 30 profession terms are female-leaning, with a mean projection of
$-1.84$ at the peak layer. We attribute this reversal to mDeBERTa's
Disentangled Attention mechanism, suggesting that architectural choices can
override data-driven gender associations. All experiments require no
fine-tuning and run on CPU using publicly released checkpoints.

\end{abstract}

% =======================================================
\section{Introduction}
% =======================================================

Modern foundation models compress linguistic diversity into unified
representational systems, embedding languages with radically different
morphologies, scripts, and training data volumes within shared parameters.
This architectural unification enables cross-lingual generalisation, yet it
also raises a structural concern: when languages with unequal resource
availability coexist within the same geometric space, how are social
associations distributed across them? If representational capacity is
unevenly shaped by pretraining data, then bias may not manifest uniformly
across linguistic contexts.

Gender bias in language models has traditionally been examined in English
through static embedding analyses. These studies demonstrate that social
associations can be recovered geometrically, suggesting that bias is encoded
in the spatial structure of representations rather than at the surface level.
Multilingual contextual models, however, introduce two additional dimensions.
First, multilingual pretraining blends languages with vastly different resource
profiles, making it unclear whether bias patterns generalise consistently.
Second, transformer representations are layered and dynamic: meaning is
progressively constructed across network depth rather than fixed at a single
stage.

A layer-wise perspective therefore becomes essential. Lower transformer layers
tend to encode lexical and morphological features, while deeper layers capture
increasingly abstract semantic relationships \citep{rogers2020primer,
tenney-etal-2019-bert}. If gender associations evolve across layers, then bias
is not merely inherited from surface statistics but may be actively reshaped
during contextual processing. Crucially, different attention architectures may
reshape this signal in fundamentally different ways---a hypothesis this paper
directly tests.

This study investigates gender bias as a joint cross-lingual, layer-wise, and
architecture-comparative phenomenon. We evaluate five multilingual transformer
models spanning three distinct vocabulary systems (WordPiece, SentencePiece
250K, SentencePiece 1M), two model scales (base, large), and two attention
mechanisms (standard multi-head, Disentangled Attention). Tigrigna---a
low-resource Ge'ez-script language---serves as our primary case study because
it provides a near-zero-shot setting: the models were not trained on Tigrigna
corpora, making any detected bias attributable to subword geometry inherited
from cross-lingual transfer rather than language-specific statistics.

\subsection{Research Questions}

\begin{itemize}[noitemsep,leftmargin=*]
  \item \textbf{RQ1}: Do multilingual transformers encode measurable gender bias
        across Arabic, Spanish, and Tigrigna under a unified geometric protocol?
  \item \textbf{RQ2}: How does measured bias vary across transformer layers, and
        at which depths do associations peak?
  \item \textbf{RQ3}: Do two independent geometric estimators agree on bias
        trends across languages and models?
  \item \textbf{RQ4}: Does attention architecture determine the direction of
        encoded gender bias, independent of pretraining data?
\end{itemize}

% =======================================================
\section{Related Work}
% =======================================================

\subsection{Bias in Static Word Embeddings}

Early work established that social biases are encoded in the geometric structure
of static word embeddings. \citet{bolukbasi2016man} showed that a gender
direction derived from anchor sets reveals systematic profession--gender
associations and proposed projection-based debiasing. \citet{caliskan2017semantics}
introduced the Word Embedding Association Test (WEAT), providing a statistical
framework for measuring implicit associations. \citet{garg2018word} demonstrated
that embedding bias correlates with historical societal trends. Later,
\citet{gonen2019lipstick} showed that removing a single bias direction does not
eliminate deeper geometric bias, motivating continued geometric analysis.

\subsection{Interpreting Contextual Representations via Reductions}

Transformer models produce contextualised representations that vary across
input sentences and network depth, complicating direct application of
geometric methods designed for static embeddings. \citet{bommasani2020reductions}
formalised a framework for interpreting pretrained contextualised representations
via reductions to static embeddings, demonstrating that contextual token
embeddings can be systematically aggregated across occurrences to produce stable
word-level vectors suitable for geometric analysis. \citet{rogers2020primer}
surveys the uneven distribution of linguistic information across transformer
layers, motivating layer-wise analysis of properties such as bias.

\subsection{Multilingual Bias and Cross-Lingual Considerations}

As multilingual transformers became widely adopted, bias analysis extended
beyond English. \citet{dev2019measuring} examined bias across languages in
multilingual embeddings, highlighting cross-lingual variation in association
strength. \citet{lauscher2020zero} studied bias transfer in shared multilingual
architectures, showing that associations from high-resource languages can
propagate across the shared parameter space.

However, existing multilingual bias studies typically evaluate only final-layer
or pooled representations, rarely investigating how bias evolves across internal
layers. Low-resource languages remain underrepresented, and no systematic
evaluation compares attention architectures under a controlled protocol.

\subsection{Architecture-Driven Effects on Representations}

\citet{he2021deberta} introduced mDeBERTa-v3, whose Disentangled Attention
separates content and positional information into independent attention matrices.
This architectural choice substantially changes how contextual relationships are
encoded across layers. Similarly, \citet{liang2023xlmv} showed that enlarging
the SentencePiece vocabulary from 250K (XLM-R) to 1M tokens (XLM-V) improves
tokenisation coverage across low-resource languages but can introduce
corpus-alignment mismatches when the expanded vocabulary fragments tokens
differently from the evaluation corpus.

\subsection{How This Work Differs}

We integrate geometric bias analysis \citep{bolukbasi2016man} with contextual-to-static
reductions \citep{bommasani2020reductions} applied layer-wise across five
multilingual architectures. The primary contribution is a controlled
architecture comparison in the near-zero-shot setting of Tigrigna, revealing
that the Disentangled Attention mechanism produces a complete reversal of gender
bias direction---a finding not previously documented in the multilingual bias
literature.

% =======================================================
\section{Methodology}
% =======================================================

\subsection{Models}

We evaluate five pretrained multilingual transformer models without fine-tuning.
All checkpoints are loaded from HuggingFace Hub; no parameter updates are
performed.

\begin{itemize}[noitemsep,leftmargin=*]
  \item \textbf{mBERT} (\texttt{bert-base-multilingual-cased},
        \citealt{devlin2019bert}): 12 transformer layers + embedding layer
        (13 hidden states), WordPiece vocabulary of $\approx$120K tokens
        covering 104 languages.
  \item \textbf{XLM-RoBERTa-base} (\texttt{xlm-roberta-base},
        \citealt{conneau2020unsupervised}): 12 layers (13 hidden
        states), SentencePiece vocabulary of 250K tokens, trained on
        CommonCrawl across 100 languages.
  \item \textbf{XLM-RoBERTa-large} (\texttt{xlm-roberta-large},
        \citealt{conneau2020unsupervised}): 24 layers (25 hidden
        states), same 250K SentencePiece vocabulary, larger model scale.
  \item \textbf{XLM-V-base} (\texttt{facebook/xlm-v-base},
        \citealt{liang2023xlmv}): 12 layers (13 hidden states),
        SentencePiece vocabulary of 1M tokens, designed to improve
        tokenisation coverage for low-resource languages.
  \item \textbf{mDeBERTa-v3-base} (\texttt{microsoft/mdeberta-v3-base},
        \citealt{he2021deberta}): 12 layers (13 hidden states),
        SentencePiece vocabulary, uses Disentangled Attention with separate
        content and position matrices---the key architectural difference
        examined in RQ4.
\end{itemize}

\subsection{Languages and Data}

We analyse three languages: Spanish (\texttt{es}), Arabic (\texttt{ar}),
and Tigrigna (\texttt{ti}).
Tigrigna is the primary case study: it is a Semitic language written in
the Ge'ez (Ethiopic) script, with negligible dedicated presence in
multilingual pretraining corpora, making it a stringent near-zero-shot
evaluation.

For each language we provide:
\begin{itemize}[noitemsep,leftmargin=*]
  \item A sentence-level corpus (one sentence per line).
  \item A list of 30 profession terms with translations.
  \item Gender anchor sets: 11 male and 10 female Tigrigna terms (pronouns,
        kinship terms, generic gender markers).
\end{itemize}

\subsection{Context Sampling}

For each target word $w$, we retrieve up to $n$ sentences containing $w$
via substring matching. We use $n=50$ contexts for profession terms and
$n=12$ for anchor terms to ensure stable centroid estimation within CPU
memory constraints. Words with fewer than the minimum available contexts
are excluded.

\subsection{Contextual-to-Static Reduction}

Transformer models produce contextualised token representations that vary
across inputs and network depth. To enable geometric analysis, we adopt the
reduction framework of \citet{bommasani2020reductions}, who demonstrate that
contextual embeddings can be systematically aggregated to produce stable
word-level vectors.

For each occurrence of word $w$ in context $c_i$, we extract hidden states
at every transformer layer. Subword token representations corresponding to
$w$ are mean-pooled across the subword fragments; the result is then averaged
across $n$ contexts. Formally, at layer $\ell$:

\begin{equation}
  \mathbf{s}_{\ell}(w)
  =
  \frac{1}{n}
  \sum_{i=1}^{n}
  \mathbf{h}_{\ell}(w, c_i),
  \label{eq:reduction}
\end{equation}

where $\mathbf{h}_{\ell}(w, c_i)$ is the mean-pooled subword representation
of $w$ at layer $\ell$ in context $c_i$ \citep{bommasani2020reductions}.
This yields one static embedding per layer for each target word.

\subsection{Gender Direction Construction}

Following \citet{bolukbasi2016man}, we define a gender direction using
predefined anchor sets. Let $\mathcal{M}$ and $\mathcal{F}$ denote the
male and female anchor term sets. For each layer $\ell$, anchor
representations are reduced as in Eq.~\eqref{eq:reduction}.
Layer-specific centroids are:

\begin{equation}
  \mathbf{c}_{M,\ell}
  =
  \frac{1}{|\mathcal{M}|}
  \sum_{w \in \mathcal{M}}
  \mathbf{s}_{\ell}(w),
  \qquad
  \mathbf{c}_{F,\ell}
  =
  \frac{1}{|\mathcal{F}|}
  \sum_{w \in \mathcal{F}}
  \mathbf{s}_{\ell}(w),
\end{equation}

and the layer-specific gender direction is:

\begin{equation}
  \mathbf{g}_{\ell}
  =
  \mathbf{c}_{M,\ell} - \mathbf{c}_{F,\ell}.
\end{equation}

This vector encodes the principal geometric contrast between male and female
anchors at layer $\ell$ \citep{bolukbasi2016man}.

\subsection{Bias Measurement}

For each profession $p$ and layer $\ell$ we compute two geometric estimators.

\paragraph{(1) Gender Direction Projection (proj).}

The scalar projection of $\mathbf{s}_{\ell}(p)$ onto the gender direction:

\begin{equation}
  \mathrm{ProjBias}_{\ell}(p)
  =
  \frac{
    \mathbf{s}_{\ell}(p) \cdot \mathbf{g}_{\ell}
  }{
    \|\mathbf{g}_{\ell}\|
  }.
  \label{eq:proj}
\end{equation}

Positive values indicate alignment toward the male centroid; negative values
indicate alignment toward the female centroid. The scalar projection preserves
vector magnitude and is sensitive to the strength of the gender association
\citep{bolukbasi2016man}.

\paragraph{(2) Centroid Cosine Difference (cosdiff).}

As a complementary magnitude-normalised estimator:

\begin{equation}
  \mathrm{CosBias}_{\ell}(p)
  =
  \cos\!\bigl(\mathbf{s}_{\ell}(p),\, \mathbf{c}_{M,\ell}\bigr)
  -
  \cos\!\bigl(\mathbf{s}_{\ell}(p),\, \mathbf{c}_{F,\ell}\bigr),
  \label{eq:cosdiff}
\end{equation}

where cosine similarity is defined as \citep{manning2008introduction}:

\begin{equation}
  \cos(\mathbf{a},\mathbf{b})
  =
  \frac{\mathbf{a} \cdot \mathbf{b}}
       {\|\mathbf{a}\|\,\|\mathbf{b}\|}.
  \label{eq:cosine}
\end{equation}

This estimator directly compares relative angular proximity to male and female
centroids and is invariant to vector scale.

\subsection{Estimator Agreement (RQ3)}

To assess measurement robustness, we compute Spearman rank correlation
between the two estimators across profession terms at each layer
\citep{spearman1904}:

\begin{equation}
  \rho_{\ell}
  =
  \mathrm{Spearman}\!\left(
    \bigl\{\mathrm{ProjBias}_{\ell}(p)\bigr\}_{p},\,
    \bigl\{\mathrm{CosBias}_{\ell}(p)\bigr\}_{p}
  \right).
  \label{eq:spearman}
\end{equation}

Spearman rank correlation measures the degree to which both estimators
rank professions in the same order, regardless of scale differences
\citep{spearman1904}. A value of $\rho_{\ell} = 1$ indicates perfect ordinal
agreement; $\rho_{\ell} = 0$ indicates no agreement.

\subsection{Layer-Wise Aggregation}

For visualisation we compute the mean projection score across retained
profession terms at each layer $\ell$, yielding a layer-wise bias trajectory:

\begin{equation}
  \overline{\mathrm{Proj}}_{\ell}
  =
  \frac{1}{|P|}
  \sum_{p \in P}
  \mathrm{ProjBias}_{\ell}(p),
\end{equation}

where $P$ is the set of retained professions. This trajectory summarises how
the overall geometric gender signal evolves across model depth.

% =======================================================
\section{Results}
% =======================================================

\subsection{Case Study: \textit{ሓኪም} (Doctor) in Tigrigna}
\label{sec:case-hakim}

\paragraph{Setup.}
We examine the Tigrigna word for \textit{doctor}, \textit{ሓኪም}, as a
representative case for the low-resource setting under XLM-RoBERTa-base.
The corpus contains 499 sentences in which \textit{ሓኪም} appears, drawn
from template-based generation covering both male and female co-occurrence
contexts (\textit{ሰበይቲ ሓኪም እያ} / \textit{ሓወይ ሓኪም እዩ}).
Gender geometry is constructed from 11 male and 10 female anchor words,
each supported by 12 corpus sentences, yielding a stable gender direction
across all 13 layers of XLM-RoBERTa-base \citep{conneau2020unsupervised}.

\paragraph{Layer-wise bias.}
Table~\ref{tab:hakim-layers} reports the layer-wise projection
(Eq.~\ref{eq:proj}) and cosine-difference (Eq.~\ref{eq:cosdiff})
estimators for \textit{ሓኪም}. The mean projection across all layers is
$+0.626$, placing \textit{ሓኪም} in the male-leaning half of the 30-profession
distribution. The signal peaks at layer~9 ($\mathrm{proj} = +2.259$),
consistent with the finding that mid-to-upper layers encode the richest
semantic associations \citep{tenney-etal-2019-bert}. Layers~2--4 yield
near-neutral projections, suggesting that the male association strengthens
as representations propagate through the feed-forward sub-layers.

\begin{table}[H]
\centering
\caption{Layer-wise bias scores for \textit{ሓኪም} (XLM-RoBERTa-base, Tigrigna).
         $\mathrm{proj}_\ell$: Eq.~\ref{eq:proj};
         $\Delta\cos_\ell$: Eq.~\ref{eq:cosdiff}.}
\label{tab:hakim-layers}
\small
\begin{tabular}{rrrrl}
\toprule
\textbf{Layer} & \textbf{proj} & \textbf{$\Delta$cos} & \textbf{Direction} \\
\midrule
 0 & $+0.374$ & $+0.014$ & male-leaning \\
 1 & $+0.250$ & $+0.013$ & male-leaning \\
 2 & $-0.082$ & $+0.007$ & female-leaning \\
 3 & $-0.017$ & $+0.006$ & neutral \\
 4 & $-0.011$ & $+0.003$ & neutral \\
 5 & $+0.614$ & $+0.001$ & male-leaning \\
 6 & $+1.296$ & $+0.002$ & male-leaning \\
 7 & $+1.577$ & $+0.001$ & male-leaning \\
 8 & $+1.195$ & $-0.002$ & male-leaning \\
 9 & $+2.259$ & $-0.005$ & male-leaning \\
10 & $+0.096$ & $-0.004$ & male-leaning \\
11 & $+0.672$ & $-0.002$ & male-leaning \\
12 & $-0.083$ & $-0.000$ & female-leaning \\
\midrule
\textbf{Mean} & $\mathbf{+0.626}$ & & \textbf{male-leaning} \\
\bottomrule
\end{tabular}
\end{table}

\paragraph{Representation uniqueness.}
To assess whether the bias signal reflects genuine contextual gender inference
or a near-static property of the token, we compute pairwise cosine similarity
across all 499 contextual vectors of \textit{ሓኪም} at each layer
(Table~\ref{tab:hakim-uniqueness}). At the final layer (layer~12) all
$\binom{499}{2} = 124{,}251$ pairs exceed the $0.99$ threshold, indicating
complete representational collapse. This is consistent with the observation
that low-resource language tokens tend to produce near-static representations
when context cannot be leveraged \citep{bommasani2020reductions}.

\begin{table}[H]
\centering
\caption{Representation uniqueness for \textit{ሓኪም} across selected layers
         ($n = 499$ sentences, $124{,}251$ unique pairs per layer).}
\label{tab:hakim-uniqueness}
\small
\begin{tabular}{rrrrrr}
\toprule
\textbf{Layer} & \textbf{mean cos} & \textbf{std} & \textbf{min} & \textbf{max} & \textbf{n\_dups ($>$0.99)} \\
\midrule
 0 & 0.983 & 0.016 & 0.951 & 1.000 &  60{,}786 \\
 1 & 0.936 & 0.038 & 0.790 & 1.000 &   4{,}145 \\
 6 & 0.951 & 0.019 & 0.882 & 1.000 &   3{,}247 \\
 9 & 0.974 & 0.011 & 0.930 & 1.000 &   6{,}363 \\
12 & 0.997 & 0.001 & 0.993 & 1.000 & 124{,}251 \\
\bottomrule
\end{tabular}
\end{table}

\paragraph{Interpretation.}
The detected male-leaning bias for \textit{ሓኪም} is real but must be
interpreted as an \emph{intrinsic token-level association} encoded during
pretraining, not as evidence of context-sensitive gender reasoning. Because
XLM-RoBERTa cannot leverage cross-sentence variation for Tigrigna, the bias
measurement functions analogously to a static embedding probe
\citep{caliskan2017semantics}---but one derived from a multilingual contextual
model rather than a frequency-based co-occurrence matrix.

% ----------------------------------------------------------
\subsection{Data Coverage After Context Filtering}
\label{sec:coverage}

Table~\ref{tab:coverage} reports vocabulary and corpus coverage after
applying the minimum context requirement.

\begin{table}[H]
\centering
\caption{Coverage after context filtering for Tigrigna (primary case study).
         Spanish and Arabic results follow the same protocol.}
\label{tab:coverage}
\small
\begin{tabular}{lccccc}
\toprule
Language & Corpus size & Professions retained & Male anchors & Female anchors & Contexts per profession \\
\midrule
Arabic   & {[..]}   & {[..]}  & {[..]} & {[..]}  & {[..]}  \\
Spanish  & {[..]}   & {[..]}  & {[..]} & {[..]}  & {[..]}  \\
Tigrigna & $\approx$14{,}850 & 30 / 30 & 11 & 10 & 50 \\
\bottomrule
\end{tabular}
\end{table}

For Tigrigna under XLM-RoBERTa-base all 30 profession terms have 498--499
corpus hits each (per representation-uniqueness analysis). Under XLM-V-base,
two terms---\textit{ጠበቓ} (lawyer) and \textit{ነዳቓይ} (mechanic)---yielded
zero corpus hits due to tokenisation mismatches introduced by the 1M
SentencePiece vocabulary (see Section~\ref{sec:xlmv}), leaving 28 terms
for that model. All other models retained the full 30 terms.

% ----------------------------------------------------------
\subsection{Five-Model Comparison on Tigrigna}
\label{sec:fivemodel}

Table~\ref{tab:model-comparison} summarises the bias profile of each
architecture on Tigrigna, constituting the main quantitative result.

\begin{table}[H]
\centering
\caption{Five-model gender bias comparison on Tigrigna (ti). Peak layer is
         the layer with the highest mean projection score (Eq.~\ref{eq:proj}).
         M:F:N = professions majority male-leaning : female-leaning : neutral
         across layers. Min~$\rho$ = minimum Spearman correlation
         (Eq.~\ref{eq:spearman}) across all layers for that model.}
\label{tab:model-comparison}
\small
\begin{tabular}{lccccccc}
\toprule
\textbf{Model} & \textbf{Layers} & \textbf{Vocab} & \textbf{Peak} & \textbf{Peak proj} & \textbf{M:F:N} & \textbf{Min $\rho$} \\
\midrule
mBERT              & 13 & WP $\approx$120K   & ---  & $-3.61^{*}$  & n/a    & $\dagger$1.000 \\
XLM-R-base         & 13 & SP 250K            &   9  & $+2.99$      & 26:3:1 & 0.992 \\
XLM-R-large        & 25 & SP 250K            &  23  & $+3.46$      & 25:3:2 & 0.992 \\
XLM-V-base         & 13 & SP 1M              &  11  & $+2.63$      & $25{:}2{:}1^{\ddagger}$ & $0.970^{\S}$ \\
mDeBERTa-v3-base   & 13 & SP $\approx$128K   &  11  & $-1.84^{\P}$ & 2:28:0 & 0.974 \\
\bottomrule
\multicolumn{7}{p{0.95\linewidth}}{%
$^{*}$~mBERT maps all Tigrigna single-word tokens to \texttt{[UNK]};
  reported value is the degenerate single-word cluster mean at layer~0.
  Multi-word professions form a second cluster: $-0.677$ at layer~0.
  No profession-level differentiation is present.\\
$^{\dagger}$~Spearman $\rho \approx 1.000$ at all layers is a tied-rank
  artifact: all professions share identical values at each layer, making
  ranks trivially tied; the result is not meaningful metric agreement.\\
$^{\ddagger}$~28/30 professions; \textit{ጠበቓ} and \textit{ነዳቓይ} excluded
  (0 corpus hits under the 1M vocabulary).\\
$^{\S}$~Layer~12 anomaly: $\rho = 0.720$; all other layers $\rho \geq 0.970$.\\
$^{\P}$~Negative peak indicates female-leaning direction; mDeBERTa produces a
  complete directional reversal relative to the XLM-R family.
} \\
\end{tabular}
\end{table}

% ----------------------------------------------------------
\subsection{Layer-Wise Bias Trajectories}
\label{sec:trajectories}

Figure~\ref{fig:projection-grid} displays the layer-wise mean projection
trajectory for all five models. The panels reveal qualitatively distinct
profiles that correspond to architectural and vocabulary differences.

\begin{figure}[H]
\centering

\begin{subfigure}[b]{0.45\linewidth}
  \includegraphics[width=\linewidth]{../output/ti/bert-base-multilingual-cased/figs/ti_mbert_projection_curve.png}
  \caption{mBERT --- WordPiece [UNK] degeneration; two clusters visible.}
  \label{fig:proj-mbert}
\end{subfigure}
\hfill
\begin{subfigure}[b]{0.45\linewidth}
  \includegraphics[width=\linewidth]{../output/ti/xlm-roberta-base/figs/ti_xlmr_base_projection_curve.png}
  \caption{XLM-RoBERTa-base --- single peak at layer~9 ($+2.99$).}
  \label{fig:proj-xlmrbase}
\end{subfigure}

\vspace{1em}

\begin{subfigure}[b]{0.45\linewidth}
  \includegraphics[width=\linewidth]{../output/ti/xlm-roberta-large/figs/ti_xlmr_large_projection_curve.png}
  \caption{XLM-RoBERTa-large --- double peak: layers~13 ($+2.99$) and~23 ($+3.46$).}
  \label{fig:proj-xlmrlarge}
\end{subfigure}
\hfill
\begin{subfigure}[b]{0.45\linewidth}
  \includegraphics[width=\linewidth]{../output/ti/facebook_xlm-v-base/figs/ti_xlmv_base_projection_curve.png}
  \caption{XLM-V-base --- peak at layer~11 ($+2.63$); collapse to $-3.27$ at layer~12.}
  \label{fig:proj-xlmv}
\end{subfigure}

\vspace{1em}

\begin{subfigure}[b]{0.45\linewidth}
  \includegraphics[width=\linewidth]{../output/ti/microsoft_mdeberta-v3-base/figs/ti_mdeberta_projection_curve.png}
  \caption{mDeBERTa-v3-base --- all layers negative from layer~1; female-leaning reversal; collapse to $-7.24$ at layer~12.}
  \label{fig:proj-mdeberta}
\end{subfigure}
\hfill
\begin{subfigure}[b]{0.45\linewidth}
  \centering
  \vspace{2em}
  \small
  \textbf{Key observations across panels:}\\[0.5em]
  \begin{tabular}{lp{0.7\linewidth}}
    \textbullet & XLM-R family (b, c, d) builds bias monotonically in mid-layers \\[0.3em]
    \textbullet & Scale amplifies: large $>$ base at peak \\[0.3em]
    \textbullet & mDeBERTa (e) is uniquely female-leaning \\[0.3em]
    \textbullet & Final layer (layer 12/24) collapses in all models \\[0.3em]
    \textbullet & mBERT (a) shows no profession differentiation \\
  \end{tabular}
  \caption{Observations.}
  \label{fig:proj-notes}
\end{subfigure}

\caption{Layer-wise mean projection bias ($\overline{\mathrm{Proj}}_\ell$) for all five
         models on Tigrigna. Positive values indicate male-leaning direction;
         negative values indicate female-leaning direction. Layers are indexed
         from the embedding layer (0) through the final transformer block.}
\label{fig:projection-grid}
\end{figure}

\FloatBarrier

\paragraph{Shared patterns.}
Four findings are consistent across the XLM-R family (XLM-R-base, XLM-R-large,
XLM-V-base):
\begin{enumerate}[noitemsep,leftmargin=*]
  \item Bias is near-zero at layer~0 and builds progressively across depth.
  \item It peaks at a middle-to-upper layer (9--11 for base-size models,
        13 and 23 for XLM-R-large).
  \item It collapses at the final layer (layer~12 for base models,
        layer~24 for the large model) toward near-zero or negative values.
  \item Increasing model scale (XLM-R-base vs.\ XLM-R-large) amplifies
        the peak magnitude ($+2.99 \to +3.46$) without changing direction.
\end{enumerate}

\paragraph{mBERT degeneration.}
mBERT's WordPiece vocabulary does not include Tigrigna characters; all tokens
are mapped to \texttt{[UNK]}. This produces two degenerate clusters at each
layer: 26 single-word professions share one identical value per layer (e.g.,
$-3.614$ at layer~0), and the 4 multi-word professions (which tokenise to
multiple \texttt{[UNK]} tokens with different position encodings) share a
distinct but also uniform value ($-0.677$ at layer~0). No profession-level
differentiation exists, confirming that mBERT constitutes a \emph{negative
control} rather than a meaningful bias measurement.

% ----------------------------------------------------------
\subsection{Most Biased Professions at Peak Layer}
\label{sec:professions}

Table~\ref{tab:top5-xlmrlarge} lists the five most male-biased Tigrigna
professions at the peak layer of XLM-RoBERTa-large (layer~23), alongside
professions with persistent female-leaning associations.

\begin{table}[H]
\centering
\caption{Most gender-biased professions at peak layer~23,
         XLM-RoBERTa-large, Tigrigna. Bias is ProjBias (Eq.~\ref{eq:proj}).}
\label{tab:top5-xlmrlarge}
\small
\begin{tabular}{llr}
\toprule
\textbf{Profession (Tigrigna)} & \textbf{English gloss} & \textbf{ProjBias$_{23}$} \\
\midrule
ዳኛ                    & judge           & $+7.13$ \\
መምህር                 & teacher         & $+6.89$ \\
ፖሊስ                  & police officer  & $+6.11$ \\
ኮኮ                   & cook/chef       & $+6.03$ \\
ሰራሕተኛ ጽሬት           & cleaning worker & $+3.93$ \\
\midrule
\multicolumn{3}{l}{\textit{Female-leaning across most layers:}} \\
ጠበቓ                  & lawyer          & (neg.\ across layers 5--19) \\
ሳይንቲስት               & scientist       & (neg.\ at early layers) \\
\bottomrule
\end{tabular}
\end{table}

Notably, \textit{ሰራሕተኛ ጽሬት} (cleaning worker) is strongly male-leaning
under XLM-R-large despite the occupation being stereotypically female in many
societies---illustrating that geometric associations may reflect pretraining
co-occurrence patterns rather than social stereotypes directly.

% ----------------------------------------------------------
\subsection{Estimator Agreement (RQ3)}
\label{sec:agreement}

Table~\ref{tab:spearman} reports Spearman rank correlation (Eq.~\ref{eq:spearman})
between the projection and cosine-difference estimators across all layers for
each model.

\begin{table}[H]
\centering
\caption{Layer-wise Spearman $\rho$ between \texttt{proj} and \texttt{cosdiff}
         estimators on Tigrigna. Minimum across all layers is reported as
         the conservative bound on agreement.}
\label{tab:spearman}
\small
\begin{tabular}{lccl}
\toprule
\textbf{Model} & \textbf{Min $\rho$} & \textbf{Max $\rho$} & \textbf{Notes} \\
\midrule
mBERT            & ---   & ---   & degenerate (tied ranks, see Table~\ref{tab:model-comparison}) \\
XLM-R-base       & 0.992 & 1.000 & perfect at layers 11--12 \\
XLM-R-large      & 0.993 & 1.000 & perfect at layers 5 and 24 \\
XLM-V-base       & $0.720^{*}$ & 0.999 & layer 12 anomaly; $\rho \geq 0.970$ for layers 0--11 \\
mDeBERTa-v3-base & 0.974 & 0.999 & lowest at layer 9 \\
\bottomrule
\multicolumn{4}{p{0.85\linewidth}}{%
$^{*}$~Layer~12 of XLM-V-base shows a representational collapse that
       disrupts both estimators simultaneously, yielding the one anomalous
       low-correlation point. Excluding layer~12, the minimum is 0.970.}\\
\end{tabular}
\end{table}

For all non-degenerate models, Spearman $\rho \geq 0.970$ at every valid layer
(addressing RQ3). The two estimators---one magnitude-sensitive
(Eq.~\ref{eq:proj}), one angle-only (Eq.~\ref{eq:cosdiff})---agree on the
ordinal ranking of professions despite their different geometric properties.
This confirms that the observed bias patterns are not artifacts of estimator
choice.

% ----------------------------------------------------------
\subsection{Architectural Reversal: mDeBERTa and Disentangled Attention}
\label{sec:reversal}

The most striking result is a complete inversion of bias direction under
mDeBERTa-v3-base. Whereas the XLM-R family produces 25--26 male-leaning
professions per model, mDeBERTa produces only 2 male-leaning professions
and 28 female-leaning ones (addressing RQ4).

Table~\ref{tab:mdeberta-exceptions} shows the only two professions that
remain male-leaning under mDeBERTa, averaged across all layers.

\begin{table}[H]
\centering
\caption{Only male-leaning professions under mDeBERTa-v3-base (Tigrigna),
         averaged across all 13 layers. All 28 remaining professions are
         female-leaning.}
\label{tab:mdeberta-exceptions}
\small
\begin{tabular}{llr}
\toprule
\textbf{Profession (Tigrigna)} & \textbf{English gloss} & \textbf{Mean ProjBias} \\
\midrule
ፕሮፌሰር & professor & $+0.27$ \\
ሓረስታይ & farmer    & $+0.13$ \\
\bottomrule
\end{tabular}
\end{table}

The layer-12 collapse is especially pronounced for mDeBERTa: individual
professions reach values as extreme as $-17.02$ (\textit{ነዳቓይ}) and
$-16.58$ (\textit{ተክኒሻን}), far exceeding the XLM-R-large layer-12 collapse
(which approaches zero, near $-0.51$ mean). This suggests that the final
hidden state in mDeBERTa undergoes a qualitatively different representational
transition than in the XLM-R family.

We attribute the directional reversal to Disentangled Attention
\citep{he2021deberta}, which encodes content and position information in
separate attention matrices. When content-based and position-based
representations are disentangled, the gender direction $\mathbf{g}_\ell$
(constructed from anchor centroids, Eq.~3) may align orthogonally or even
inversely with profession subspace geometry relative to standard multi-head
attention. The result is that profession vectors cluster on the female side
of the gender axis---not because mDeBERTa associates professions with
femininity in any direct cultural sense, but because its internal geometry
is organised differently.

% ----------------------------------------------------------
\subsection{Vocabulary Size and Corpus Alignment: XLM-V}
\label{sec:xlmv}

XLM-V-base was designed to overcome the vocabulary bottleneck in XLM-R by
expanding the SentencePiece vocabulary from 250K to 1M tokens
\citep{liang2023xlmv}. For most low-resource languages this improves
tokenisation fidelity. However, for Tigrigna we observe a paradoxical
misalignment: two profession terms---\textit{ጠበቓ} (lawyer) and
\textit{ነዳቓይ} (mechanic)---are tokenised differently by the 1M vocabulary
than by the 250K vocabulary, producing subword sequences absent from the
evaluation corpus. Both terms receive 0 corpus hits under XLM-V, and are
excluded from the bias calculation.

This illustrates that larger vocabulary coverage does not guarantee
better corpus alignment: if the corpus was generated or tokenised with a
different vocabulary, enlarging the model vocabulary introduces new
mismatches. The XLM-V results (28 professions, peak $+2.63$) are otherwise
consistent with the XLM-R family, confirming that when tokenisation succeeds,
bias patterns are robust to vocabulary size.

% =======================================================
\section{Discussion}
% =======================================================

\subsection{Cross-Lingual Patterns}

The Tigrigna results demonstrate that multilingual transformers encode
gender associations even for languages absent from their dedicated
pretraining data. This is consistent with \citet{lauscher2020zero}, who
show that properties from high-resource languages transfer implicitly to
low-resource ones via shared subword representations. Because Tigrigna
tokens are processed via SentencePiece subword pieces that partially overlap
with other Semitic or phonologically similar languages, the models inherit
geometric associations from those languages' training signal.

However, the nature of this transferred bias is not transparent. The
associations detected for Tigrigna professions cannot be verified against
ground-truth human judgments of Tigrigna gender norms, and may reflect
the gender coding of phonologically similar tokens in higher-resource
languages. This highlights a critical limitation of cross-lingual
geometric bias analysis: the source of the detected signal is opaque.

\subsection{Layer-Wise Interpretation}

The consistent pattern across the XLM-R family---near-zero at layer~0,
building through mid-layers, peaking at layer~9--11 (base) or layer~23
(large), then collapsing---aligns with the established view that lower
layers encode syntactic and morphological features while upper layers
encode semantic relationships \citep{tenney-etal-2019-bert,rogers2020primer}.
Gender associations, as social-semantic properties, are expected to peak
in the layers where semantic organisation is most pronounced.

The final-layer collapse is notable: bias scores for most professions
approach zero or reverse sign at layer~12 (base) or layer~24 (large).
This suggests that the model's final-layer representations are pulled
toward a language-model prediction objective that partially averages out
profession-specific gender associations. Practitioners using final-layer
embeddings for downstream tasks may therefore underestimate the degree
of bias present in intermediate representations.

\subsection{Architecture Determines Bias Direction (RQ4)}

The mDeBERTa reversal is the paper's most important finding. mBERT,
XLM-R-base, XLM-R-large, and XLM-V-base all encode male-leaning bias
for the same Tigrigna professions trained on broadly comparable multilingual
corpora. mDeBERTa-v3-base, trained on similar multilingual data, produces
female-leaning associations for 28 of the same 30 professions.

The only systematic difference between mDeBERTa and the XLM-R family is
the attention mechanism: Disentangled Attention versus standard multi-head
attention. This is strong evidence that \emph{architectural choice
determines bias direction} in the low-resource cross-lingual setting,
overriding the signal from pretraining data. This has practical implications:
bias mitigation strategies that focus on pretraining data curation may be
insufficient if the attention mechanism itself reshapes geometric associations.

\subsection{Metric Robustness}

The agreement between the projection estimator (Eq.~\ref{eq:proj}) and the
cosine-difference estimator (Eq.~\ref{eq:cosdiff}) at $\rho \geq 0.970$
across all valid models is a methodological positive. The two estimators
capture different geometric properties---one is magnitude-sensitive, the
other is angle-only---yet they rank professions in nearly identical order
at every layer. This confirms that the bias ranking is a robust geometric
feature of the model's representation space, not an artifact of the
particular estimator chosen.

The one exception is XLM-V layer~12, where $\rho = 0.720$. Inspection
of the data shows that the representation collapse at this layer
simultaneously disrupts both estimators in different ways, breaking their
rank agreement. This is consistent with a catastrophic representational
transition at the final layer, not with genuine metric disagreement.

\subsection{Limitations}

\begin{itemize}[noitemsep,leftmargin=*]
  \item \textbf{Corpus construction}: The Tigrigna corpus is template-based,
        controlling for surface-level context but limiting ecological validity.
        Bias measurements may differ on naturally occurring text.
  \item \textbf{Tokenisation and morphology}: Tigrigna is morphologically rich;
        profession terms may be systematically fragmented across models,
        affecting the subword mean-pooling step.
  \item \textbf{Anchor coverage}: The 21-word anchor set (11 male, 10 female)
        is necessarily limited. Different anchor choices may shift the
        gender direction and alter profession rankings.
  \item \textbf{Near-static representations}: As shown in
        Section~\ref{sec:case-hakim}, Tigrigna tokens produce highly similar
        representations across contexts. The bias signal reflects intrinsic
        token geometry rather than contextually modulated gender inference.
  \item \textbf{No human validation}: We have no ground-truth Tigrigna
        speaker judgments about profession gender associations to validate
        against, limiting the cultural interpretation of findings.
  \item \textbf{Causal opacity}: We cannot directly observe the mechanism by
        which Disentangled Attention reverses bias direction; the attribution
        to attention architecture is inferential.
\end{itemize}

% =======================================================
\section{Conclusion}
% =======================================================

This study presents the first systematic layer-wise gender bias analysis
of five multilingual transformer architectures applied to Tigrigna, a
low-resource Ge'ez-script language. Using a unified geometric protocol
grounded in contextual-to-static reductions \citep{bommasani2020reductions}
and projection-based bias estimation \citep{bolukbasi2016man}, we draw four
conclusions.

First, four of five models encode measurable male-leaning gender associations
for Tigrigna professions despite the language's absence from dedicated
pretraining data, confirming cross-lingual transfer of social associations
(RQ1). Second, bias accumulates progressively across transformer depth,
peaking at mid-to-late layers and collapsing at the final layer---a pattern
consistent with the semantic organisation of transformer representations
(RQ2). Third, two independent geometric estimators agree at Spearman
$\rho \geq 0.970$ across all valid models, confirming methodological
robustness (RQ3). Fourth, and most importantly, mDeBERTa-v3-base's
Disentangled Attention mechanism completely inverts the bias direction
relative to the XLM-R family, producing female-leaning associations for
28 of 30 professions---strong evidence that architectural choice can
override data-driven gender associations (RQ4).

These findings have concrete implications for multilingual NLP practitioners:
(1) geometric bias analysis requires SentencePiece vocabulary coverage as a
prerequisite for low-resource languages---WordPiece models like mBERT are
degenerate negative controls; (2) bias audit results depend critically on
which model layer is inspected; (3) bias mitigation strategies focusing
solely on pretraining data may be insufficient if the attention mechanism
independently reshapes bias geometry.

\section{Future Work}

\begin{itemize}[noitemsep,leftmargin=*]
  \item Extend the analysis to Arabic and Spanish layer-wise trajectories
        under all five models for a complete cross-lingual comparison.
  \item Collect small-scale human judgments from Tigrigna speakers to
        validate profession-gender associations against cultural norms.
  \item Probe the Disentangled Attention mechanism mechanistically
        (attention-head attribution, representation similarity analysis)
        to understand why it reverses bias direction.
  \item Test whether post-hoc debiasing methods \citep{bolukbasi2016man}
        applied at the intermediate peak layer reduce downstream bias in
        classification tasks.
  \item Evaluate additional low-resource languages with Ge'ez and other
        non-Latin scripts to assess the generalisability of the findings.
\end{itemize}

% =======================================================
\clearpage
\bibliographystyle{plainnat}
\bibliography{refs}

% =======================================================
\appendix
\section{Profession Inventory}

\subsection{Full Tigrigna Profession List (30 Terms)}

\begin{itemize}[noitemsep, leftmargin=*, columns=2]
  \item ሓኪም (doctor)
  \item ነርስ (nurse)
  \item መምህር (teacher)
  \item ፕሮፌሰር (professor)
  \item መሃንዲስ (engineer)
  \item ሳይንቲስት (scientist)
  \item ጠበቓ (lawyer)
  \item ዳኛ (judge)
  \item ሓላፊ (manager)
  \item ዳይረክተር (director)
  \item ወተሃደር (soldier)
  \item ፖሊስ (police officer)
  \item ሓረስታይ (farmer)
  \item ነጋዳይ (merchant)
  \item መካኒክ (mechanic)
  \item ኤለክትሪሻን (electrician)
  \item ፈለኛሞ (pilot)
  \item ነዳቓይ (miner/driller)
  \item ኣርኪቴክት (architect)
  \item ኣካውንታንት (accountant)
  \item ባንከር (banker)
  \item ስነ ጥበባዊ (artist)
  \item ሙዚቀኛ (musician)
  \item ተዋሳኢ (actor)
  \item ስራሕ ፈጣሪ (entrepreneur)
  \item ኮኮ (cook/chef)
  \item ሰራሕተኛ ጽሬት (cleaning worker)
  \item ሓላዊ ጸጥታ (security guard)
  \item ተክኒሻን (technician)
  \item ፕሮግራመር (programmer)
\end{itemize}

\subsection{Sample Gender Anchors (Tigrigna)}

\textbf{Male anchors (11):} pronouns, kinship terms, and generic masculine
markers in Tigrigna (e.g., \textit{ንሱ} `he', \textit{ሓወይ} `my brother',
\textit{ኣቦ} `father', \textit{ወዲ} `son/boy').

\textbf{Female anchors (10):} corresponding feminine forms
(e.g., \textit{ንሳ} `she', \textit{ሓፍተይ} `my sister',
\textit{ኣደ} `mother', \textit{ጓል} `daughter/girl').

\end{document}
